\Cref{tab:2-position} relates the preceding methods to JHQ-GPU.
Small-table methods motivate exploiting code structure; graph and
binary-code indexes provide alternative speed--capacity trade-offs;
adaptive-effort methods provide the quality-control comparison.
Chapters 4--6 develop these design choices, and Chapter 7 evaluates them.

\begin{table}[!htbp]
\centering
\caption[Related approaches and JHQ-GPU design questions]{Related approaches and the design questions they motivate for JHQ-GPU.}
\label{tab:2-position}
\footnotesize\setlength{\tabcolsep}{4pt}
\begin{tabularx}{\linewidth}{@{}>{\raggedright\arraybackslash}p{.20\linewidth}*{2}{>{\raggedright\arraybackslash}X}@{}}
\toprule
Approach & Main trade-off & Question for JHQ-GPU \\
\midrule
IVF-PQ / OPQ & Learned codebooks reduce distortion but need not separate coordinates internally & How can the primary code's coordinate structure simplify GPU scoring? \\
Fast Scan / Quicker ADC & Small tables accelerate lookup under hardware and precision constraints & Does splitting tables repay extra GPU lookups? \\
SONG / CAGRA & Graph navigation reduces scoring but adds graph storage and dependent access & Where can a compressed scan compete with resident graph search? \\
BANG & Host storage extends capacity at the cost of communication & What can a complete compressed index resident on one GPU achieve? \\
IVF-RaBitQ & Code size and probe effort balance distance-estimation quality and speed & When does residual refinement justify its extra work? \\
Adaptive effort & Quality control may require labels, prediction or calibration & Can workload agreement select a useful budget at acceptable quality and cost? \\
\bottomrule
\end{tabularx}
\end{table}
